\subsection{Package Hallucination}
% Foundational and survey work has examined hallucination in natural language generation and general LLM settings, including dialogue, summarization, and factual question answering~\cite{ji2023survey,huang2025survey,zhang2025siren,alansari2026large}. Recent software engineering studies show that hallucination also appears in code generation, where models may produce incorrect APIs, non-existent libraries, or flawed implementation logic~\cite{liu2024exploring,tian2025codehalu,zhang2025llm}; documentation-grounded generation has also been studied as a mitigation for API hallucinations~\cite{jain2025mitigating}.
% Package hallucination is a more specific failure mode in which an LLM recommends dependency names that do not exist or are not valid for the target ecosystem. Spracklen et al.~\cite{spracklen2025we} conducted a large-scale empirical study of package hallucinations, evaluating 16 code LLMs on 576,000 generated code samples and showing that the problem is pervasive across models and programming languages. Krishna et al.~\cite{krishna2025importing} further showed that package hallucination rates vary across programming languages, model families, model scales, and task specificity, and remain non-negligible even for stronger coding models. HFuzzer~\cite{zhao2025hfuzzer} complements these empirical studies by using phrase-based fuzzing to trigger package hallucinations in dependency-related tasks.
% Package hallucination also creates software supply-chain risk because developers or coding agents may pass a fabricated name to a package manager. Prior work on open-source supply chains has studied package-confusion and related registry attacks~\cite{ladisa2023sok,neupane2023beyond,kaplan2021survey}, while package-hallucination studies and industry reports show that hallucinated names can become attractive targets for slop-squatting or malicious registration~\cite{spracklen2025we,krishna2025importing,lanyado2024diving}.

Prior work has studied hallucination in natural language generation and general LLM settings, such as dialogue, summarization, and factual question answering~\cite{ji2023survey,huang2025survey,zhang2025siren,alansari2026large}. 
Recent software engineering studies show that hallucination also appears in code generation, where models may produce incorrect APIs, non-existent libraries, or flawed implementation logic~\cite{liu2024exploring,tian2025codehalu,zhang2025llm}; documentation grounding has also been explored to mitigate API hallucinations~\cite{jain2025mitigating}. 
Package hallucination is a specific failure mode where an LLM recommends dependency names that do not exist or are unavailable in the target ecosystem. 
Empirical studies show that package hallucination is common across models, programming languages, model sizes, and task settings~\cite{spracklen2025we,krishna2025importing}, and HFuzzer further triggers such failures through phrase-based fuzzing~\cite{zhao2025hfuzzer}. 
This problem also raises software supply chain risks, since developers or coding agents may pass fabricated names to package managers, creating opportunities for package confusion, slop squatting, or malicious registration~\cite{ladisa2023sok,neupane2023beyond,kaplan2021survey,lanyado2024diving,spracklen2025we,krishna2025importing}.

Researchers have proposed several strategies to mitigate hallucinations~\cite{tonmoy2024comprehensive,alansari2026large}. 
Retrieval methods ground generation with external knowledge sources, self-correction prompts models to verify and refine their outputs, and model adaptation changes model behavior through fine-tuning or similar training procedures~\cite{jain2025mitigating,ji2023towards,madaan2023self,spracklen2025we}. 
Inference-time and post-hoc editing methods, such as TruthX and PURR, mitigate hallucinations by intervening in internal representations or revising generated outputs without full retraining~\cite{zhang2024truthx,chen2023purrefficientlyeditinglanguage}. 
However, these methods often depend on external resources, require repeated inference, revise outputs without correcting the model's default behavior, or incur high training costs. 
In contrast, we address package hallucination through localized model editing with a small set of editing examples.


% To mitigate hallucinations, researchers have proposed several strategies~\cite{tonmoy2024comprehensive,alansari2026large}. Retrieval-based methods augment generation with external knowledge sources to improve factual grounding, including API-documentation grounding for code generation~\cite{jain2025mitigating}. Self-correction approaches prompt models to verify and refine their own outputs through iterative reasoning and self-reflection~\cite{ji2023towards,madaan2023self}.
% Model-adaptation methods directly modify model behavior through supervised fine-tuning or similar training procedures~\cite{spracklen2025we}.
% Inference-time and post-hoc editing methods, such as TruthX and PURR, mitigate hallucinations by intervening in internal representations or revising generated outputs without full retraining~\cite{zhang2024truthx,chen2023purrefficientlyeditinglanguage}.
% While these techniques can reduce hallucination rates, they often depend on external knowledge sources, require repeated inference, filter outputs without correcting the model's default behavior, or incur substantial retraining costs. Different from these generation-time pipelines or full adaptation procedures, we address package hallucinations through localized model editing with a small set of editing examples.

\subsection{Model Editing}
% Model editing, also known as knowledge editing, aims to update specific knowledge or behaviors in a pretrained model without full retraining~\cite{sinitsin2020editable,zhu2020modifying,wang2024knowledgeediting}. Its central goal is to correct targeted model errors while preserving the model's original capabilities. Prior work commonly evaluates editing methods along reliability, generalization, locality, and portability~\cite{zhang2024comprehensive,wang2024easyedit}. These criteria are particularly important in software engineering tasks, where an edit should fix a target error without degrading the model's broader coding ability. Existing methods generally follow two paradigms. Weight-preserved methods, such as SERAC, GRACE, and IKE, rely on external memories or in-context demonstrations to override selected predictions without directly changing the original model parameters~\cite{mitchell2022memory,hartvigsen2023aging,zheng2023edit}. In contrast, weight-modified methods directly update model parameters through learned editors or localized parameter changes. For example, KnowledgeEditor and MEND learn auxiliary networks to generate targeted parameter updates~\cite{de2021editing,mitchell2022mend}, while later studies show that aggressive or sequential editing may harm general abilities or cause catastrophic forgetting~\cite{gu2024modeleditingharms,gupta2024catastrophic}.

% Among weight-modified methods, localized editing methods are especially relevant to our work. Motivated by evidence that knowledge is encoded in specific transformer components~\cite{geva2021transformer,dai2022knowledge}, ROME, MEMIT, and PMET edit selected internal modules to rewrite factual associations at different scales~\cite{meng2022rome,meng2023memit,li2024pmet}. Recent studies further extend model editing beyond factual triples, including detoxification, code LLM robustness enhancement, and editing for LLMs4Code~\cite{wang2024detoxifying,liu2025creme,li2025model}. However, package hallucination remains unexplored in model editing. Unlike factual updates that rewrite isolated subject-object associations, package hallucination requires correcting a task-conditioned validity boundary over package-like names. Our work targets this package-validity boundary editing problem.

Model editing, also known as knowledge editing, aims to correct specific knowledge or behaviors in pretrained models without full retraining~\cite{sinitsin2020editable,zhu2020modifying,wang2024detoxifying}. 
A successful edit should fix the target error while preserving the model's original capabilities, and is commonly evaluated by reliability, generalization, locality, and portability~\cite{zhang2024comprehensive,wang2024easyedit}. 
Existing methods can be broadly divided into two groups. 
Weight-preserving methods, such as SERAC, GRACE, and IKE, rely on external memories or in-context demonstrations to override selected predictions without modifying model parameters~\cite{mitchell2022memory,hartvigsen2023aging,zheng2023edit}. 
Weight-modified methods directly update model parameters through learned editors or localized changes. 
For example, KnowledgeEditor and MEND learn auxiliary networks to generate parameter updates~\cite{de2021editing,mitchell2022mend}, while later studies show that aggressive or sequential editing may harm general abilities or cause forgetting~\cite{gu2024modeleditingharms,gupta2024catastrophic}.

Among weight-modified methods, localized editing methods are especially relevant to our work.
Motivated by evidence that knowledge is encoded in specific transformer components~\cite{geva2021transformer,dai2022knowledge}, ROME, MEMIT, and PMET edit selected internal modules to rewrite factual associations at different scales~\cite{meng2022rome,meng2023memit,li2024pmet}. 
Recent work further extends model editing beyond factual triples, including detoxification, code LLM robustness, and editing for LLMs4Code~\cite{wang2024detoxifying,liu2025creme,li2025model}. 
However, package hallucination remains unexplored in model editing. 
Unlike factual updates that rewrite isolated subject-object associations, package hallucination requires correcting a task-conditioned package-validity boundary over package names. 
Our work targets this package-validity boundary editing problem.
